\section{Related Work}
\label{sec:related_work}

Automatic hate speech and offensive language detection has been studied extensively through shared tasks and benchmark datasets. OffensEval introduced a hierarchical view of offensive language identification, distinguishing between offensive and non-offensive posts, targeted and untargeted offenses, and the type of target \citep{zampieri2019offenseval}. Its multilingual continuation showed that the task is not limited to English and that cross-lingual evaluation remains difficult \citep{zampieri2020offenseval}. HatEval focused on hate speech against women and immigrants, making the distinction between generic offensive language and identity-targeted abuse especially relevant \citep{basile2019hateval}. These tasks motivate our five-class setup, where insults are separated from racism, homophobia, and sexism.

The paper by \citet{akhtar2021whose}, included in the project topic, is particularly relevant because it frames hate speech annotation as a perspective problem. The authors argue that abusive language is subjective and that annotators with different cultural, demographic, or social backgrounds may disagree on whether a message is hateful, offensive, aggressive, or stereotypical. Instead of treating disagreement only as noise, they view it as useful information about how different communities perceive online abuse.

Their work introduces a perspective-aware approach: annotators are grouped by similar perspectives, separate gold standards are created, and supervised models are trained on group-specific labels. The authors also propose an inclusive ensemble model that marks an instance as positive if at least one perspective-aware classifier detects abuse. Their multi-perspective BREXIT Twitter corpus includes four related categories: hate speech, aggressiveness, offensiveness, and stereotype.

This perspective is useful for our project even though our Romanian datasets do not contain full annotator-level information. Our schema keeps an \texttt{annotation\_quality\_flag}, separates generic insults from identity-targeted categories, and treats racism, homophobia, and sexism as distinct classes. The work of \citet{akhtar2021whose} also supports our decision not to automatically label keyword-based candidates as hate speech, since identity-group mentions can be neutral, quoted, anti-discriminatory, or abusive depending on context.
